\documentclass[12pt,a4paper]{article}

% For pdfLaTeX: uses palatino package
% For LuaLaTeX/XeLaTeX: comment out \usepackage{palatino} and uncomment:
%   \usepackage{fontspec}
%   \setmainfont{TeX Gyre Pagella}  % Palatino-compatible OpenType font
\usepackage{palatino}
\usepackage{amsmath,amssymb,amsthm}
\usepackage{geometry}
\usepackage{microtype}
\usepackage{hyperref}
\usepackage{setspace}
\usepackage{csquotes}
\usepackage{titlesec}
\usepackage{abstract}
\usepackage{booktabs}
\usepackage{array}

\geometry{
  top=1.25in,
  bottom=1.25in,
  left=1.25in,
  right=1.25in
}

\onehalfspacing

\hypersetup{
  colorlinks=true,
  linkcolor=black,
  citecolor=black,
  urlcolor=black
}

\titleformat{\section}{\large\bfseries}{\thesection.}{0.75em}{}
\titleformat{\subsection}{\normalsize\bfseries}{\thesubsection.}{0.75em}{}

\theoremstyle{plain}
\newtheorem{theorem}{Theorem}
\newtheorem{proposition}{Proposition}

\theoremstyle{definition}
\newtheorem{definition}{Definition}
\newtheorem{hypothesis}{Hypothesis}
\newtheorem{principle}{Principle}

\theoremstyle{remark}
\newtheorem{remark}{Remark}

\title{\textbf{Motor Phonology and Symbolic Reachability:\\
Keyboards, Scripts, and the Generative Minimum}}

\author{Flyxion}

\date{June 2026}

\begin{document}

\maketitle

\begin{abstract}
\noindent
Keyboard layouts are conventionally understood as spatial maps: twenty-six
physical locations, one per letter, organized to minimize finger travel or
optimize frequency-weighted input. This paper argues that the conventional
understanding is too shallow and proposes an alternative theoretical framework
in which keyboard layouts are better understood as \emph{motor phonologies}---
structured systems that partition a continuous space of bodily movements into
stable symbolic categories, in a manner directly analogous to the way natural
languages partition acoustic or articulatory space into phonemes.

The argument proceeds through three connected claims. First, skilled typists
represent letters not primarily as spatial addresses but as motor categories
organized around finger identity and directional trajectory. The finger is the
place of articulation; the direction is the manner. Second, learning a new
keyboard layout is therefore not merely learning new locations but acquiring
a new motor phonology: carving a distinct set of contrastive categories from
the space of possible hand movements. Third, the design of keyboard interfaces
can be understood as a problem of \emph{symbolic reachability}: given a fixed
motor vocabulary, how large and how structured can the space of accessible
distinctions be made?

These three claims are developed through an analysis of a family of keyboard
projects---the Eight-Letter Keyboard, and a set of script-projection keyboards
mapping the learned QWERTY motor geometry onto Arabic, Phoenician, Unicode
Italic, and Standard Galactic scripts---that collectively instantiate a
recurring design principle: preserve the generative path, change the
destination. The paper situates this principle within a broader theoretical
framework connecting motor compression, distributed cognition, and
accessibility geometry, and proposes that the same question underlies
efforts ranging from the design of minimal input devices to the long-term
preservation of symbolic civilizations.
\end{abstract}

\bigskip

\section{The Phonological Analogy}

The standard account of keyboard learning proceeds as follows. A novice
approaches the keyboard as a spatial puzzle: twenty-six letters occupy
twenty-six positions, and the task is to memorize the spatial address of each.
Practice converts explicit spatial lookup into automatized motor execution,
and expertise is achieved when the typist no longer needs to consciously
retrieve locations. The underlying representation, on this account, is a
lookup table that has been made automatic: one spatial address per symbol,
stored as procedural memory.

This account is likely wrong, or at least incomplete, in a way that has
significant consequences for how we understand both keyboard design and motor
learning more generally.

Consider the analogous account of phonological acquisition in natural language.
A child learning a language does not memorize a table of phonemes and then
automatize its retrieval. She undergoes a prolonged process of categorical
perception formation in which a continuous acoustic space---variations in
formant frequency, voice onset time, place of articulation, and dozens of
other parameters---becomes structured into a discrete set of contrastive
categories. The category boundary between \textsc{p} and \textsc{b} in English,
for instance, falls at a particular voice onset time that English speakers
have internalized as a categorical distinction, while speakers of other languages
may draw the boundary differently or not at all. What is learned in phonological
acquisition is not a table but a \emph{partition}: a structured carving of a
continuous perceptual-articulatory space into regions that function as discrete
symbolic units.

The proposal of this paper is that keyboard learning involves an analogous
process. The space being partitioned is not acoustic but motor: it is the space
of possible hand movements, characterized by dimensions including which finger
acts, whether the movement is upward, central, or downward, how much force is
applied, and what temporal rhythm the stroke occupies. Keyboard learning is the
process by which this continuous motor space becomes partitioned into discrete
symbolic categories. Expertise is not the automatization of a lookup table but
the stabilization of a motor phonology: a system of contrastive motor categories
that generates symbolic output through their differentiation.

This analogy is more than suggestive. It has direct experimental implications,
connects keyboard design to a rich theoretical tradition in linguistics and
cognitive science, and yields a significantly different conception of what
good keyboard design should optimize.

\section{Finger Identity as Place of Articulation}

\subsection{The Articulatory Parallel}

In the phonology of consonants, two primary dimensions characterize the
articulatory gesture: the place of articulation (where in the vocal tract
the constriction occurs) and the manner of articulation (what kind of
constriction it is). Bilabial stops are produced at the lips with a complete
closure. Alveolar fricatives are produced at the alveolar ridge with a partial
constriction that produces turbulence. These two dimensions, together with
voicing, generate most of the consonant inventory of the world's languages from
a relatively small number of categorical values: roughly five or six places of
articulation, six or seven manners, and a voicing binary.

The Eight-Letter Keyboard proposes that the motor phonology of typing has
precisely this structure. The finger is the analogue of place of articulation:
it specifies which articulator is engaged. The directional parameter---whether
the stroke is toward the upper, home, or lower position in the finger's
channel---is the analogue of manner: it specifies how the articulator moves.
Together, these two parameters generate the full alphabetic inventory through
a coordinate system of the form
\[
  \text{Symbol} \;=\; \text{Finger} \;\times\; \text{Direction},
\]
in direct analogy to the articulatory coordinate system
\[
  \text{Phoneme} \;=\; \text{Place} \;\times\; \text{Manner} \;\times\; \text{Voicing}.
\]

\subsection{Evidence for Finger Identity as Cognitive Primitive}

The phonological analogy would be merely formal if there were no evidence that
finger identity functions as a genuine cognitive primitive in motor memory.
Several lines of evidence suggest that it does.

Experienced touch typists, when asked to recall properties of individual
letters, reliably exhibit an asymmetry: they can identify the finger that
types a given letter more accurately and more rapidly than they can identify
its absolute row position on the keyboard. Many typists who are entirely
unable to point to the key for a given letter can correctly identify whether
it is a left-hand or right-hand key and which finger of that hand is responsible.
This pattern is exactly what the phonological analogy predicts: the place
of articulation---which articulator is engaged---is more robustly stored than
the precise spatial position of the gesture.

Further evidence comes from interference and transfer experiments in piano
playing, which offers a closely related motor phonology. Pianists show
finger-specific interference effects: a passage trained with one finger
transfers differently to the anatomically adjacent finger than to a non-adjacent
finger on the same hand, and differently again to the homologous finger on the
opposite hand. These transfer patterns are inconsistent with a purely spatial
representation of key positions and consistent with a representational scheme
organized around finger identity, with other parameters specified relative to
that anchor.

The analogy also predicts a specific pattern of relearning difficulty across
different types of keyboard change. Switching to a layout that preserves all
finger assignments but alters row positions---the directional parameter---should
be substantially easier than switching to a layout that scrambles finger
assignments, because the former requires relearning only the secondary parameter
while the latter requires relearning the primary categorical structure. This
prediction distinguishes the motor phonology hypothesis from the spatial address
hypothesis, which predicts approximately equal difficulty for the two types of
change, and it is directly testable.

\subsection{Switching Layouts as Phonological Restructuring}

The motor phonology framework reframes the experience of switching keyboard
layouts in a way that matches the phenomenological reports of people who have
undergone the transition. Typists who switch from QWERTY to Dvorak or Colemak
consistently describe the experience not as learning where new keys are but as
a period of perceptual and motor confusion followed by a gradual stabilization
of a new movement grammar. The letters lose their locations and then acquire
new ones, but what is rebuilt is not a table but a feel: a sense of rightness
and wrongness about the movement that operates below the level of conscious
spatial retrieval.

This description fits poorly with the lookup table account, which would predict
that the transition involves learning twenty-six new locations and that the
difficulty would be roughly proportional to the number of locations changed.
It fits well with the phonological account, which predicts that the transition
involves restructuring a set of motor categories and that the difficulty reflects
the degree of categorical reorganization required. Complete reorganization---
moving from QWERTY to a layout with different finger assignments---would be a
full motor phonology change, comparable to a native speaker acquiring the
phonological system of a typologically unrelated language. Partial
reorganization---a layout that preserves finger assignments while altering
directional parameters---would be more like a dialect shift: familiar in its
primary categories but requiring adjustment in secondary features.

This is precisely the logic underlying the Eight-Letter Keyboard. If finger
identity is the primary categorical dimension of motor phonology, then a
keyboard organized explicitly around that dimension should be the most learnable
for typists who already possess the finger-identity structure, because it
asks them to relearn only the directional parameter while preserving the primary
categorical architecture of their existing motor phonology.

\section{Script Projection and the Preservation of Motor Paths}

\subsection{A Family of Reachability Tools}

A family of keyboard projects makes the motor phonology framework concrete from
a different direction. Rather than reorganizing the alphabet into a more
cognitively transparent structure, these projects address the problem of
reaching entirely different symbolic systems---Arabic, Phoenician, Unicode
Italic, Standard Galactic---without requiring the user to abandon the motor
program acquired through years of QWERTY practice.

Each project defines a mapping
\[
  f : L \;\to\; S,
\]
where $L$ is the Latin alphabet as typed on a standard QWERTY keyboard and $S$
is the target script. The mapping is constructed to be as mnemonically
transparent as possible: phonetically related sounds are assigned to the same
key, historically related letterforms are placed at corresponding positions, and
the shift key functions as a semantic modifier that accesses the more emphatic
or deeper articulatory variant of each sound.

In the Arabic keyboard, for example, the mapping preserves the phonetic
relationship between Latin and Arabic consonants wherever possible:
\texttt{s} types \textit{sin} (the Arabic sibilant \textsc{s}), \texttt{S}
types \textit{sad} (the emphatic counterpart, produced with pharyngealization),
\texttt{h} types \textit{ha} (the ordinary glottal fricative), and \texttt{H}
types \textit{hah} (the pharyngeal fricative, a sound without a close Latin
equivalent). The design allows a user
with Latin keyboard motor memory to begin producing Arabic text immediately,
learning the Arabic letter-sound correspondences through the act of typing
rather than through a separate memorization exercise.

The Phoenician keyboard follows the historical etymology of the Latin alphabet
itself. The Latin letter \texttt{a} descends from Phoenician \textit{aleph};
\texttt{b} from \textit{beth}; \texttt{d} from \textit{daleth}; \texttt{k}
from \textit{kaph}; \texttt{m} from \textit{mem}; and so on. The Phoenician
keyboard formalizes this etymology as a keyboard mapping, so that typing in
Phoenician using the layout is simultaneously tracing the historical ancestry
of the Latin letters being pressed.

The Unicode Italic keyboard is structurally the simplest: it maps each Latin
letter to its Unicode mathematical italic variant (e.g., \texttt{a}
$\to$ $\mathit{a}$, \texttt{b} $\to$ $\mathit{b}$), leaving the alphabetic
structure entirely unchanged and altering only the typographic rendering. The
interest of this case lies precisely in its simplicity: it shows that the
operation the other keyboards perform---mapping a motor geometry onto a distinct
symbol space---is already taking place in any keyboard layout, because the
relationship between a keystroke and a specific Unicode codepoint is always a
projection rather than an identity.

The Standard Galactic keyboard takes the furthest step. Standard Galactic
Alphabet is an invented script, developed and extended over many years, with no
historical connection to Latin or Arabic. The keyboard maps each QWERTY key to
its Standard Galactic equivalent, making an entirely novel symbolic system
accessible through an already-mastered motor program. The user reaches a
genuinely different distinction space---a script with its own visual identity,
its own letterform logic, its own cultural associations---while performing
essentially the same motor acts as QWERTY typing.

\subsection{The Common Structure}

Despite their apparent variety, all four projects instantiate the same
underlying operation. The motor geometry of QWERTY typing---the set of learned
finger movements, their associated categories, their kinematic patterns---is
treated as a fixed infrastructure. The symbolic output of that infrastructure
is then changed by interposing a projection layer between the motor act and
the rendered glyph. The typist does not learn a new motor program; she learns
a new interpretation of her existing motor program.

\begin{principle}[Preserve the Path, Change the Destination]
A symbolic domain $S$ becomes reachable from an existing motor vocabulary $M$
through a projection $f : M \to S$ that is learned at the level of symbolic
interpretation rather than motor reorganization. The cognitive cost of
reaching $S$ is then determined primarily by the complexity of $f$, not
by the cost of relearning $M$.
\end{principle}

This principle has a direct connection to the admissibility framework developed
elsewhere. The motor geometry $M$ represents an already-acquired reachability
structure: a set of accessible symbolic positions reachable through learned
motor acts. The projection $f$ extends this reachability structure to a new
symbolic domain without destroying the existing one. An admissible such
projection is one that preserves enough of the structure of $M$ that the user's
existing motor knowledge can guide the acquisition of $f$---that the mnemonic
bridges embedded in the projection allow the new symbolic domain to be reached
from the familiar motor ground.

The Arabic and Phoenician keyboards are more admissible in this sense than an
arbitrary key-scramble would be, because their projections are structured around
phonetic and etymological relationships that create real mnemonic bridges.
The Standard Galactic keyboard is less mnemonically structured, which means its
admissibility depends more heavily on practice and less on prior knowledge.
The Unicode Italic keyboard is perfectly admissible in the limit: the projection
is bijective on the letter level and changes only the glyph rendering, so no
new symbolic knowledge is required at all.

\section{Hierarchy and Generative Depth}

\subsection{Complexity Through Structure, Not Surface}

A conventional keyboard grows in complexity as symbol inventories grow: more
symbols require more keys. This is the flat expansion model, in which the
complexity of the interface scales linearly with the complexity of the symbol
space it must access. The motor phonology framework suggests an alternative:
complexity can grow through \emph{depth} rather than surface area, through
hierarchical structure rather than proliferation of independent elements.

The prefix-key systems of modal editors provide the clearest existing
demonstration of this principle. In editors organized around hierarchical key
sequences---Vim's command mode, Spacemacs's leader-key system, the GNU Emacs
key binding architecture---a small number of physical keys generates access
to an enormous command vocabulary. A single key establishes a namespace; a
second key selects within that namespace; further keys specify within sub-
namespaces. The structure is a tree rather than a flat lookup, and the depth
of the tree determines the size of the accessible vocabulary without requiring
any growth in the number of physical inputs.

The formal structure is
\[
  \text{Command} \;=\; \text{Prefix} \;\times\; \text{Suffix}
\]
for two-level systems, or more generally
\[
  \text{Command} \;=\; k_1 \;\times\; k_2 \;\times\; \cdots \;\times\; k_n
\]
for $n$-level hierarchies, where each $k_i$ is drawn from a small set of
primitive inputs. This is structurally identical to the motor phonology
coordinate system $\text{Symbol} = \text{Finger} \times \text{Direction}$,
and to the articulatory coordinate system $\text{Phoneme} = \text{Place} \times
\text{Manner} \times \text{Voicing}$. All three are instances of the same
abstract principle: generate a large inventory of distinct symbols from a small
inventory of primitive dimensions through Cartesian combination.

\subsection{The Generative Minimum}

The deeper question raised by this family of design problems is what might be
called the generative minimum: the smallest set of motor primitives capable of
generating the full symbolic inventory required by a given domain.

For the Latin alphabet, the Eight-Letter Keyboard proposes a specific answer:
eight finger identities combined with three directional parameters per channel
(six for the index fingers) suffice to generate all twenty-six letters plus
additional positions for high-frequency punctuation. The generative minimum for
the alphabet is, on this account, the eight-dimensional motor primitive space
of the human hand.

This framing connects the keyboard design problem to a much older question
in the theory of symbolic systems: how much complexity must a generative
substrate contain in order to reconstruct a given symbolic inventory? The
genetic code answers this question for protein synthesis: four nucleotides,
organized into triplet codons, generate sixty-four codons, of which sixty-one
code for twenty amino acids. The phonological systems of human languages answer
it for spoken language: roughly thirty to forty phonemes, organized into syllabic
structures, generate the lexical inventories of natural languages. The
combinatorial explosion from a small generator to a large inventory is the
characteristic signature of generative systems.

\begin{definition}[Generative Minimum]
Given a target symbol space $S$ and a class of motor primitives $\mathcal{P}$,
the generative minimum is the smallest set $G \subseteq \mathcal{P}$ such that
every element of $S$ is reachable as a combination of elements of $G$ under
an admissible combinatorial rule $\rho : G^n \to S$.
\end{definition}

The Eight-Letter Keyboard proposes that the generative minimum for the Latin
alphabet, relative to the motor primitive space of the human hand, is the set
of eight finger identities. The combinatorial rule is the specification of a
directional parameter. The claim is not merely practical but cognitive: that
this particular generative structure is the one that best matches the
architecture of human motor memory and therefore the one that is most
efficiently learned, retained, and executed.

\section{Preservation, Archives, and the Reed Wall}

\subsection{Symbolic Reachability as a Preservation Problem}

The keyboard design problems discussed in this paper can be understood as special
cases of a more general preservation problem. Symbolic systems---scripts,
languages, command vocabularies, musical notations, mathematical symbolisms---
represent accumulated distinctions that have been found useful or meaningful
by the communities that developed them. Their value lies in the distinctions
they make available: the difference between Arabic emphatic and non-emphatic
consonants, between Phoenician letters and their descendants, between Standard
Galactic glyphs and the Latin letters that map to them. These distinctions
constitute intellectual infrastructure.

The problem of symbolic reachability is the problem of making these distinctions
accessible to users who did not grow up with the native motor program for
producing them. A native Arabic speaker learns the motor phonology of Arabic
writing directly; the acquired gesture for \textit{sad} (the emphatic sibilant)
is as natural as the gesture for \textsc{s} in English. A person who already
commands the Latin motor
phonology and wishes to access Arabic faces a different challenge: not the
absence of phonological capacity but the absence of the specific motor-symbolic
coordinate system that Arabic writing requires.

The script-projection keyboards address this challenge by interposing a
projection layer that makes Arabic, Phoenician, or Standard Galactic reachable
through an already-acquired motor geometry. The approach is analogous to the
construction of a translation lexicon that allows a speaker of one language to
access the distinctions available in another without requiring full
acquisition of the target language's phonological system. The distinctions
remain genuine---the Arabic emphatic consonants are genuinely different from
their non-emphatic counterparts---but the path to them has been routed through
a familiar motor infrastructure.

\subsection{The Ark Principle}

There is a deeper structural parallel between the keyboard design projects and
the problem of preserving complex symbolic systems across time and societal
disruption. The intuition is ancient: the Ark was never interpreted, in most
theological traditions, as carrying every possible thing. It carried enough: a
breeding population, a seed stock, a generative minimum capable of
reconstructing what would otherwise be lost. The principle is preservation
through compression, and it recurs in every domain where the problem of
carrying forward complex information through a bottleneck arises.

Libraries are not full replicas of the worlds they document; they are curated
selections whose archival value lies in the distinctions they preserve and in
the accessibility of those distinctions to later users. Genetic sequences
preserve biological diversity not by carrying fully formed organisms but by
encoding the generative instructions from which organisms can be reconstructed.
Memory palaces preserve large inventories of recalled items not by storing
them verbatim but by encoding them as navigable structures from which they can
be reconstructed through traversal.

The keyboard projects instantiate this principle at the scale of individual
motor programs. The QWERTY motor geometry is a small learned structure---a
generative minimum for the Latin alphabet---from which a much larger space of
symbolic distinctions becomes accessible through projection layers. The Arabic,
Phoenician, Italic, and Standard Galactic keyboards each extend this
reachability to a new symbolic domain without requiring the user to start from
nothing. The existing motor geometry functions as an ark: it carries enough
structure to reach the new symbolic world.

\begin{principle}[Generative Preservation]
A symbolic system $S$ is preserved through a bottleneck $B$ if and only if
the generative minimum $G \subseteq B$ necessary to reconstruct $S$ survives
the bottleneck. Preservation of the full inventory of $S$ is not required;
preservation of the generative structure is sufficient.
\end{principle}

This principle applies to the keyboard design problems as a design criterion:
a good motor-phonology keyboard is one whose generative minimum is
small enough to be learned, stable enough to survive the transition from
novice to expert, and expressive enough to make the full target symbolic
inventory accessible through admissible combination.

\section{Distributed Cognition and the Embodied Symbolic System}

\subsection{Fingers as Specialized Agents}

The motor phonology framework suggests a natural interpretation of the
skilled typist's cognitive situation in terms of distributed cognition.
In the Eight-Letter Keyboard, each finger constitutes a semi-autonomous
channel with its own symbolic domain: the left little finger has access
to one region of the alphabet, the left ring finger to another, the index
fingers to larger central regions. A word emerges from the coordinated
activation of multiple channels rather than from the sequential retrieval
of independent spatial addresses.

This is not merely a metaphor. Motor neuroscience consistently finds that
skilled manual sequences involve distributed processing across cortical and
subcortical structures, with finger-specific motor programs operating with
considerable independence before being integrated into larger action sequences.
The independence of finger representations in motor cortex is well established;
somatotopic organization means that the representations of adjacent fingers
are anatomically adjacent, producing the observed pattern of differential
transfer between fingers of different proximity.

On an extended cognition view, the keyboard itself participates in this
distributed processing. The physical affordances of the keyboard---the
tactile feedback, the key spacing, the resistance profile---are part of the
cognitive system that produces text, not merely a transduction medium that
converts internal symbolic representations into physical signals. A keyboard
designed around the motor phonology of the hand is therefore a cognitive tool
designed to mesh with the distributed structure of manual motor memory, rather
than a spatial lookup table that the hand must learn to navigate.

\subsection{Movement Before Symbol}

There is a developmental dimension to the motor phonology framework that connects
it to broader theories of symbolic acquisition. The thesis that symbolic
complexity emerges from motor primitives through structured combination is
consistent with developmental accounts in which movement, proprioception, and
bodily self-differentiation precede and ground the acquisition of explicit
symbolic systems. A child learns the distinctions of its body---direction,
force, position, rhythm, bilateral asymmetry---before acquiring the linguistic
symbols that describe those distinctions. The body is the first distinction
machine.

The developmental evidence is precise on the relevant sequence. Between
approximately ten and fourteen months, infants develop deictic gestures---
pointing, showing, reaching---that function as requests or references to objects
before spoken words serve the same function. Crucially, gesture-symbol
combinations appear before symbol-symbol combinations. The child first
establishes a communicative act through motor action, and the symbol is learned
into that already-functional gestural slot. This is not merely an interesting
developmental sequence; it establishes a logical priority. The gesture is not
derived from or dependent on the symbolic output. The symbolic output is learned
into the gesture.

The pattern extends across species. Language-enculturated great apes trained to
use lexigrams or sign language reliably use reaching and showing gestures to
indicate intent before using the corresponding symbolic tokens for the same
purpose, and continue to use gestural reference in contexts where symbolic access
is uncertain. This cross-species commonality suggests that gestural communication
is not a human developmental convenience but a phylogenetically earlier and more
fundamental mode of intentional reference. Symbol use is a later, higher-cost
layer built over a gestural substrate that is both evolutionarily older and
cognitively more basic.

The keyboard, on this account, is a prosthesis for the body's distinction-making
capacity: an external structure that extends the range of distinctions the body
can generate and record. A motor-phonology keyboard makes this prosthetic
relationship explicit by organizing the symbolic space around the same
categorical dimensions---which articulator, what manner of movement---that
the body already uses to organize motor action.

The sequence in the development of the skilled typist mirrors, in compressed
form, the sequence in the development of the language-using child: first
movement, then category, then symbol. The novice learns to move fingers to
locations; the intermediate learner develops finger-level categories; the expert
operates with a fully stabilized motor phonology in which symbols are generated
directly from categorical motor acts without the intermediate step of spatial
lookup. The parallels with phonological acquisition---from continuous acoustic
variation through categorical perception to fluent speech production---are
structural rather than merely analogical.

\section{The Gesture as Primary Object}

\subsection{What Current Keyboard Systems Discard}

The motor phonology framework, developed so far at the level of individual
finger categories and their directional parameters, points toward a further and
more radical thesis. The keyboard does not merely map letters to locations.
The keyboard is a systematic instrument for discarding information.

A skilled typist does not produce a sequence of keystrokes. She produces a
coordinated, overlapping, temporally extended gesture distributed across both
hands. Multiple fingers are active simultaneously; contact is sustained across
varying durations; trajectories converge and diverge; and reversals of direction
under sustained contact introduce structured inflection points that divide the
gesture stream into phrase-like units. These features are not incidental. They
are stable, reproducible aspects of trained motor behavior. They are, however,
almost entirely discarded by existing input systems.

Standard keyboard pipelines perform a systematic flattening. Continuous and
concurrent motor activity is reduced to discrete key events. Overlapping states
are serialized by scan order. Temporal structure is compressed into ordering.
Inter-hand coordination is rendered invisible entirely. What remains is a
symbolic stream that is easy to process downstream but no longer reflects the
relational geometry of the gesture that produced it.

Three specific dimensions of loss deserve explicit attention. The first is
simultaneity. Keystrokes within a window of tens of milliseconds are
perceptually and motorically co-produced, forming what are effectively
micro-chords. Standard keyboard controllers serialize these into a total order
determined by scan timing rather than any property of the gesture. The partial
ordering of simultaneously active fingers is destroyed. The second is duration
and release structure. Every key has a corresponding release, and the interval
between onset and release encodes a second temporal signal largely independent
of ordering. Current text-entry systems ignore this almost entirely: keyup events
are captured at the hardware level but rarely used by higher-level software,
which treats text entry as defined exclusively by keydown order. The third and
most consequential loss is inter-hand coordination. The two hands form a coupled
system with stable, learnable patterns: convergence, divergence, asymmetric
anchoring, and synchronization. These patterns align with higher-level structure
in both language and music, yet current input systems provide no representation
of them.

The spacebar provides a concrete illustration of the cost of discarding release
structure. In a gesture-aware system, word boundaries emerge naturally from the
moment at which all fingers disengage simultaneously: a full release event that
is unambiguous, physically real, and produced without any deliberate additional
act. The spacebar exists precisely because the serial symbol model cannot use
this signal. Having discarded release structure, the system has no mechanism for
detecting word boundaries and must reintroduce them as an explicit symbol. The
spacebar is not a fundamental input primitive. It is an artifact of lossy
compression applied upstream.

\subsection{The Canonical Gesture}

The alternative to treating the keyboard as a sequence generator is to treat it
as a gesture sensor whose primary object is a temporally extended, relational
motor act. On this account, the symbol is not the primary object. The symbol
is a projection of the gesture---a downstream derived quantity---and the gesture
is what the motor system actually produces and what skilled motor memory actually
stores.

To make this precise, it is useful to introduce the notion of a canonical
gesture. Let the eight non-thumb fingers define eight channels, each operating
over a small constrained domain. At any moment, each channel has an activation
state, a position within its domain, and a direction of motion. The full system
state at time $t$ is the aggregation of these eight channel states. A gesture
is a time-indexed trajectory through this state space over the interval from
initial contact to full release.

The canonical gesture is obtained from the raw event stream by a
canonicalization operator that suppresses timing jitter while preserving three
topological features. The first is overlap: which channels were simultaneously
active. The second is pivot: where reversals of direction occurred under
sustained contact, marking the most constrained moments of the gesture and the
natural boundaries of phrase-like units. The third is segmentation: where full
release divided one gesture from the next, providing the word boundary that the
spacebar currently approximates.

These three features---overlap, pivot, and segmentation---are the stable
vocabulary that motor memory learns. They are invariant under the spatial
distortion of the keyboard, temporal jitter, and changes to the symbolic
interpretation layer. Two distinct raw event streams that share the same
overlap pattern, pivot structure, and segmentation points are realizations of
the same canonical gesture and will feel like the same word to the typist who
produces them.

The canonical gesture decomposes naturally into pivot-centered units: a
convergence phase, a pivot, and a divergence phase, which together constitute
what might be called a slingshot arc. The pivot is the moment of maximal
constraint---where one or more fingers reverse direction while remaining in
contact. It corresponds structurally to the syllabic peak in speech and the
harmonic anchor in music. The word \textit{machine}, for example, decomposes
into three such arcs: the approach and reversal at \textit{m}--\textit{a},
the spread to the \textit{ch} cluster and reversal at \textit{h}, and the
return through \textit{i}--\textit{n}--\textit{e} terminating in full release.
No spacebar is required; the word boundary emerges from the release structure
of the gesture itself.

\subsection{Coarticulation and the Limits of Discrete Classification}

The canonical gesture framework reveals a structural inadequacy in any approach
that treats symbolic output as the primitive object of keyboard input. A
skilled typist does not produce the letter \textit{e} and then independently
produce the letter \textit{d} and then independently produce the letter
\textit{i}. Each finger movement is shaped by anticipatory preparation for
subsequent movements and residual dynamics from prior movements. This is
coarticulation: the same property that makes phonemes in speech context-sensitive
deformations of one another rather than invariant acoustic objects.

The implications for recognition and design are parallel to those in speech.
In fluent speech, phoneme tokens do not exist as acoustically invariant objects.
They are continuous, overlapping, mutually deforming trajectories in articulatory
space. Listeners do not decode speech by isolating acoustically invariant
phoneme tokens; they reconstruct the latent articulatory trajectory responsible
for the continuously evolving acoustic field. The same logic applies to typing.
The gesture for a familiar word is not the sum of independently executed letters.
It is a globally optimized trajectory under motor and linguistic constraints,
in which the identity of any local segment depends on its position within the
larger arc.

This means that the symbolic output is semantically underdetermined from any
single local observation. The correct interpretation of a local finger
trajectory depends on what precedes and follows it. Systems that classify
keystrokes in isolation---treating each keydown event as an independent unit
to be matched against a predefined vocabulary---are modeling the wrong
ontology. The gesture is the primary object; the discrete symbol is a derived
abstraction over that gesture; and recognition requires integrating evidence
across the full temporal arc rather than committing to interpretations at
each local instant.

\subsection{Projection Functors and the Admissibility Layer}

The canonical gesture framework supports a precise account of how symbolic
output relates to the underlying motor act. A projection functor maps canonical
gesture space into an output domain: text, musical pitch-time structure, visual
form, or any other symbolic modality. The key property of such a functor is
that it depends only on the canonical gesture, not on the raw event stream. Two
raw event streams that canonicalize to the same gesture produce the same output
under any projection.

This immediately yields the design principle behind the script-projection
keyboard family. The same canonical gesture space---the QWERTY motor
phonology---is mapped into different output domains by interposing different
projection functions. The Arabic keyboard, the Phoenician keyboard, the Standard
Galactic keyboard: each is a different projection functor from the same
canonical gesture space. The motor program is invariant; only the interpretation
changes.

The projection functor account also yields a natural account of the admissibility
layer: the constraint structure that governs which projections are coherent.
A projection is admissible when the canonical gesture closes into a valid element
of the output domain---a complete word, a well-formed harmonic phrase, an
unambiguous command. Admissibility is not a binary predicate but a graded
quantity: a gesture may be more or less admissible under a given projection
depending on how completely and unambiguously it maps to a valid output.

The obstruction to admissibility---the failure of a gesture to close coherently
under projection---corresponds geometrically to the failure of local gesture
arcs to glue into a globally consistent output. If two adjacent arcs in the
canonical gesture decompose into local projections that cannot be consistently
assembled into a single output element, the obstruction is not a local error but
a global structural incompatibility. The gesture has internally consistent parts
that cannot be combined into a coherent whole under the chosen projection. This
is the keyboard-input analogue of the user-independent recognition problem in
gesture recognition: locally consistent individual realizations that fail to
assemble into a globally coherent semantic object.

The theoretical consequence is that the keyboard, properly understood, is not
a character-entry device. It is a gesture sensor equipped with a configurable
projection layer and an admissibility constraint that determines which
gesture-projection pairs produce stable, valid symbolic output. The problem
of keyboard design is not the problem of assigning characters to spatial
positions. It is the problem of choosing projection functors that are maximally
admissible with respect to the motor phonology of the users who must operate
them.

\section{From Keyboards to Reachability Geometry}

\subsection{The Generalization}

The keyboard design projects discussed in this paper are instances of a more
general problem structure. In each case, there is a space of possible
distinctions $S$ (the full symbolic inventory of a script or command language),
a motor primitive space $M$ (the set of accessible finger movements), and a
projection $f : M \to S$ that determines which distinctions are reachable and
at what cost. The design problem is to choose $f$ so as to maximize the
admissibility of the mapping---to ensure that the structure of $M$ is preserved
as fully as possible in the structure of $f(M)$---while making as much of $S$
accessible as possible from the resources that $M$ provides.

This is a reachability problem, and it connects the keyboard design projects
to a broader family of problems in which a cognitive or physical agent must
maximize the symbolic or operational distinctions accessible to it given
the constraints of its available primitives. The human hand, with its eight
non-thumb fingers and their associated directional capabilities, is a
finite motor system with a specific architecture. The question is not how
many symbols can be assigned to keys in a spatial layout but how much
symbolic space can be made reachable through the generative structure that
the hand's architecture makes available.

\subsection{Depth Instead of Surface}

The answer that the motor phonology framework suggests is consistent across
all the cases examined: the accessible distinction space can be made much
larger than the naive account would suggest, provided that the mapping between
motor primitives and symbols is organized hierarchically rather than flatly.
A flat mapping---one symbol per motor position, with no combinatorial
structure---limits the accessible vocabulary to the number of
distinguishable positions. A hierarchical mapping---symbols generated as
combinations of primitive dimensions---scales combinatorially with the number
of dimensions and the number of categorical values per dimension.

Human hands have eight non-thumb fingers, each capable of at least three
categorical directional positions. A flat mapping generates at most
$8 \times 3 = 24$ symbols---not quite enough for the Latin alphabet.
A two-level hierarchical mapping---in which the finger and direction are
each a categorical dimension of a symbol coordinate---generates exactly
the same 24 symbols, but does so through a structured generative system
that matches the cognitive architecture of motor memory rather than
requiring independent retrieval of 24 distinct spatial addresses.

The deeper insight is that the complexity of the accessible distinction space
should be engineered into the \emph{structure} of the mapping rather than
the \emph{surface} of the interface. More symbols do not require more keys;
they require deeper structure in the relationship between motor acts and
their symbolic interpretations. This is the principle that prefix-key systems
like Spacemacs have discovered operationally and that the motor phonology
framework articulates theoretically.

\begin{principle}[Depth Principle]
The complexity of the symbolic distinction space accessible through a motor
interface grows with the depth of the hierarchical structure of the
motor-to-symbol mapping, not with the size of the physical interface surface.
\end{principle}

\subsection{The Smallest Sufficient Generator}

The question that underlies all the projects discussed in this paper is, in its
most abstract form: what is the smallest structure that can survive and still
regenerate the larger system?

This question is ancient in its mythological and practical expressions. The
ark carries a breeding population rather than every individual organism. The
library preserves the structure of knowledge rather than every instantiated
text. The genetic code preserves the generative instructions for life rather
than the full inventory of organisms. In each case, preservation through a
bottleneck is possible because the system possesses generative structure: a
small generator that, when placed in the right conditions and supplied with
the right resources, can reconstruct what would otherwise be lost.

The Eight-Letter Keyboard asks whether the Latin alphabet can survive inside
eight fingers. The script-projection keyboards ask whether Arabic, Phoenician,
and Standard Galactic can be reached from the motor program for QWERTY. Both
questions are instances of the same deeper inquiry: how much symbolic
complexity is preserved in the generative minimum, and what are the conditions
under which that minimum suffices to reconstruct the full inventory?

The answer, across all the cases, appears to be: more than naive inventory
counting would suggest, provided the generative structure is organized around
the architecture of the system that must use it. A motor phonology organized
around the natural categorical structure of manual gesture can reach more of
the symbolic inventory than a flat spatial address system organized around
the physical geography of a keyboard surface, because it aligns the generative
structure with the cognitive structure of the system that must generate from it.

That alignment---between the structure of the generator and the structure of
the cognitive architecture that operates it---is, in the end, the design
principle that the motor phonology framework is trying to articulate. And it
is a principle that extends well beyond keyboards, into every domain where
the problem of preserving and accessing complex symbolic systems from a
finite generative base arises.

\section{The Binary Decision Tree and Trajectory Attractors}

\subsection{Nested Dichotomies in Manual Motor Space}

The coordinate system $\text{Symbol} = \text{Finger} \times \text{Direction}$
captures the structure of the Eight-Letter Keyboard at the level of final
categorical dimensions, but it understates the hierarchical depth of the
decision process that produces a symbol. When the anatomy of the arrangement
is examined closely, the keyboard is not organized as a flat Cartesian product.
It is organized as a binary decision tree, in which each node poses a
left-versus-right or up-versus-down question, and the leaf nodes are the
individual symbols.

The tree has the following structure. The first distinction is the most global:
which arm initiates the gesture, left or right. This is the coarsest semantic
partition, and it is not arbitrary. Left-hand letters and right-hand letters
are spatially and semantically distributed in a way that roughly mirrors the
left-right hemispheric organization of language processing, and the two hands
in skilled typing are genuinely semi-independent agents whose coordination
produces words rather than whose competition selects letters.

Within each arm the next distinction is which side of the hand: the outer
fingers (little and ring) versus the inner fingers (middle and index). This
is again a spatial and functional distinction, not merely a positional one.
The outer fingers tend to handle less frequent letters and boundary functions;
the inner fingers handle higher-frequency letters and the richer directional
vocabulary of the index column.

Within each finger the next distinction is vertical: up, home, or down. This
maps onto the three keyboard rows, but it is more naturally understood as a
directional motor parameter---the direction of the keystroke relative to the
resting position---than as a location in a spatial grid. The finger does not
navigate to a row; it extends upward, rests at home, or curls downward.

The index fingers add one further distinction: the inward column. Because the
index finger is longer and more mobile than the others, it commands two vertical
columns rather than one. The distinction between the outer and inner index
columns is a lateral movement---a drawing of the finger toward the midline of
the body---superimposed on the vertical parameter. This is not a third
categorical dimension of the same kind as the others; it is an additional
binary choice available only to the index fingers because of their superior
range of motion.

The full decision tree can be written schematically as follows, where each
node represents a binary question and each terminal node is a symbol:

\begin{center}
\textit{Left or Right?} \\[4pt]
$\hookrightarrow$ \textit{Outer or Inner fingers?} \\[4pt]
$\hookrightarrow$ \textit{Which finger?} \\[4pt]
$\hookrightarrow$ \textit{Up, Home, or Down?} \\[4pt]
$\hookrightarrow$ [\textit{Index only:} Inner or Outer column?] \\[4pt]
$\hookrightarrow$ \textbf{Symbol}
\end{center}

This tree structure has a striking property: the number of binary decisions
required to reach any leaf is small and approximately equal across all leaves.
Five decisions suffice to specify any letter in the alphabet, and the decisions
are organized in decreasing order of spatial scale: from arm, to hand side,
to finger, to row, to column. This is precisely the hierarchical structure that
the nervous system is known to prefer for motor control.

The pinkies occupy a special position in this tree. The left pinky, in addition
to its letter assignments (\texttt{Q}, \texttt{A}, \texttt{Z}), is the natural
locus for the meta-symbolic control functions: Shift, Ctrl, Alt, Win, Tab,
and Escape. The right pinky handles Enter, Backspace, Delete, and punctuation
closures. This is not a coincidence of keyboard geography. The pinkies are the
boundary fingers, the least recruited in high-frequency letter typing, and the
natural candidates for functions that modify or frame the symbolic content
produced by the other fingers. In linguistic terms, the pinkies are behaving
more like grammar than vocabulary: they modulate the interpretation of
everything else rather than contributing content themselves. The keyboard
therefore already encodes a content-versus-control distinction in the anatomy
of hand assignment, with the inner fingers carrying the lexical load and the
outer fingers managing the meta-symbolic framing.

\begin{proposition}[Binary Depth of the Alphabet]
The Latin alphabet of twenty-six letters is encodable in a binary decision
tree of depth at most five, using the anatomical distinctions of arm, hand
side, finger, vertical direction, and lateral column.
\end{proposition}

\begin{proof}
The tree yields $2 \times 2 \times 2 \times 3 \times 2 = 48$ potential
terminal nodes for the index fingers, reduced to $6 \times 3 = 18$ for the
remaining fingers per hand, giving a total capacity of $30$ for the
eight-finger assignment as shown in the Generative Sufficiency theorem.
Since $30 \geq 26$, all twenty-six letters fit within the tree. The maximum
depth from root to any leaf is five decisions: arm, hand side, finger, direction,
and (for index fingers) column.
\end{proof}

\subsection{Symbols as Trajectory Attractors}

The binary tree description captures the \emph{categorical} structure of the
motor phonology, but it still treats the selection of a symbol as the outcome
of a discrete decision process. The phenomenology of expert typing suggests
something richer: for a trained typist, the sequence of decisions collapses
into something more like navigation through a landscape than traversal of a
tree. This transition has a precise mathematical description.

Let $\mathcal{M}$ denote the continuous space of possible hand configurations
and trajectories. A typing event is not a point in $\mathcal{M}$ but a
trajectory $\gamma : [0, T] \rightarrow \mathcal{M}$. Two trajectories are
symbolically equivalent if they produce the same output:
\[
  \gamma_1 \sim \gamma_2
  \quad\Longleftrightarrow\quad
  \sigma(\gamma_1) = \sigma(\gamma_2).
\]
The alphabet is then a quotient of the trajectory space:
\[
  \Sigma \;=\; \Gamma / {\sim},
\]
where $\Gamma$ is the space of all motor trajectories. A symbol $s$ is not a
movement. It is the set of movements that preserve the same symbolic distinction:
\[
  s \;=\; [\gamma] \;=\; \{\gamma' : \sigma(\gamma') = \sigma(\gamma)\}.
\]

This formulation is stronger than the flat coordinate description because it
naturally accounts for motor variation. No two executions of the letter
\texttt{E} are physically identical. What remains invariant across executions
is not the trajectory but membership in an equivalence class. The symbol is
the class, not any particular member.

In dynamical terms, each equivalence class $[\gamma]$ corresponds to an
attractor basin in motor space. A skilled typist has learned to navigate toward
these attractor basins reliably and efficiently. The training process is not the
memorization of target positions but the carving of a landscape with stable
attractor regions---one per symbol---separated by barriers that the typist
learns to avoid.

\begin{definition}[Symbol Volume]
Let $d_{\mathcal{M}}$ be a metric on motor trajectory space. The volume of a
symbol $s$ is
\[
  V(s) \;=\; \mathrm{Vol}([\gamma_s]),
\]
the measure of the motor equivalence class of $s$ in $\mathcal{M}$.
\end{definition}

\begin{theorem}[Symbol Stability]
The robustness of a symbol $s$ under motor variation is monotonically
increasing in $V(s)$. A symbol occupying a large attractor basin is more
easily learned, more resistant to execution error, and more tolerant of
inter-user motor variation than a symbol occupying a narrow basin.
\end{theorem}

\begin{proof}
A motor perturbation $\epsilon$ is tolerated without symbolic error if and only
if the perturbed trajectory $\gamma + \epsilon$ remains within $[\gamma_s]$.
The probability of tolerance is therefore proportional to the fraction of the
$\epsilon$-ball around $\gamma$ that lies within $[\gamma_s]$, which is
monotonically increasing in $V(s)$ for any rotationally symmetric distribution
of perturbations.
\end{proof}

This theorem transforms keyboard design into a geometric optimization problem.
The ideal motor-phonology keyboard maximizes the volume of each symbol's
attractor basin subject to the constraint that basins are disjoint:
\[
  \text{maximize} \quad \sum_{s \in \Sigma} V(s)
  \quad\text{subject to}\quad
  V(s_i) \cap V(s_j) = \varnothing \;\text{ for } i \neq j.
\]
The design question is no longer which spatial location to assign to each letter
but how to arrange the motor-phonological coordinate system so that every letter
occupies the largest possible, most clearly separated attractor basin consistent
with the anatomy of the hand.

\subsection{From Coordinates to Trajectories: The Compression of Expertise}

The binary tree and attractor landscape accounts describe the structure that
expertise is working toward, but the trajectory of learning from novice to
expert reveals something further. The transition is not merely from slow to
fast execution of the same process. It is a qualitative change in the
representational format that underlies typing.

A novice traverses the decision tree explicitly: left or right arm, which
finger, which direction, and so on. Each letter requires a separate descent
through the tree, and a word requires as many descents as it has letters. The
cognitive cost of typing a word is proportional to its letter count.

An intermediate learner has automatized the individual descents but still
retrieves letters sequentially. The tree traversal has become implicit, but
each leaf is still reached independently.

An expert typist does neither. Common words have been compiled into single
motor objects: the word \textit{therefore} is not eight sequential letter
retrievals but a unified trajectory $\gamma_{\textit{therefore}}$ whose shape
the motor system has learned to execute as a whole. This is the trajectory
compression that chunking theory predicts and phenomenological reports confirm:
expert typists report that familiar words feel like single gestures, not
sequences of decisions.

\begin{hypothesis}[Trajectory Compression]
Let $w = (\ell_1, \ell_2, \ldots, \ell_n)$ be a word. The cognitive cost
of typing $w$ satisfies
\[
  C(w)
  \;\approx\;
  \begin{cases}
    \displaystyle\sum_{i=1}^n C(\ell_i) & \text{(novice)}, \\[8pt]
    C(\gamma_w) & \text{(expert)},
  \end{cases}
\]
where $\gamma_w$ is the single compiled motor trajectory for the word. Expert
cost is approximately independent of word length for frequent words, because
the trajectory is retrieved as a unit rather than assembled from parts.
\end{hypothesis}

This is the keyboard analogue of the transition from note-by-note reading to
phrase-level musical execution that musicians report in the acquisition of
fluency. On the piano, individual notes eventually give way to shapes,
intervals, and gesture trajectories: a scale is not eight separate notes but
a path with a characteristic kinematic signature. On the Eight-Letter Keyboard,
a word eventually becomes a motor melody: the individual finger-direction
coordinates are waypoints along a trajectory whose overall shape is the
primary object of retrieval.

The analogy to Swype-style gesture keyboards is exact and illuminating. A Swype
user begins by explicitly tracing through letter positions on a virtual grid.
With practice, the grid recedes and the word becomes a single continuous gesture
whose shape is recognized holistically. The individual letter positions are no
longer the representational primitives; the trajectory is. The same transition
is available on the Eight-Letter Keyboard: the binary decisions (left or right,
up or down) are scaffolding for the learning process, but what gets stored after
sufficient practice is the compiled motor trajectory for the whole word.

This convergence with Swype is not coincidental. Both systems are organizing
symbolic access around the geometry of continuous motor trajectories rather than
the enumeration of discrete positions. The Eight-Letter Keyboard makes the
initial learning more transparent by exposing the binary decision structure; Swype
makes the expert representation more visible by removing the discrete key
affordances entirely. Both arrive at the same endpoint: a navigable motor
manifold whose stable regions are words rather than letters.

\subsection{Motor Space as a Place-Cell Map}

The most radical formulation of the trajectory attractor view connects the
keyboard to the spatial navigation literature. Place cells in the hippocampus
fire when an animal is in a specific location in space, and their collective
activity constitutes a cognitive map of the environment: a landscape of
activation rather than a list of coordinates. The navigator does not remember
a sequence of positions; the navigator learns a landscape and can infer current
position and plan trajectories through that landscape.

Expert typing may be organized by an analogous cognitive structure. The skilled
typist does not remember a lookup table of letter-to-location assignments.
The typist has learned a motor landscape whose stable regions---the attractor
basins---are the symbols. Navigation through this landscape is the act of
typing. Words are not sequences of symbol retrievals but paths through the
landscape, and the motor system plans and executes these paths as continuous
trajectories from an initial state to a terminal release, just as the spatial
navigator plans a route from origin to destination.

This framing has a testable prediction. If typing is landscape navigation
rather than lookup retrieval, then the errors typists make should not be
uniformly distributed across the alphabet. They should cluster near the
boundaries between adjacent attractor basins---the saddle points of the motor
landscape---just as navigation errors in spatial environments cluster near
the boundaries between adjacent place-field regions. Substitution errors in
typing should preferentially involve letters whose motor-phonological coordinates
are adjacent in the binary tree (same arm, same hand side, same finger, adjacent
direction) rather than letters that are merely spatially proximate on the
keyboard surface. This prediction distinguishes the attractor landscape account
from the spatial address account and is in principle testable with existing
keystroke dynamics data.

The deepest consequence of the landscape formulation is its reversal of the
assumed direction of representation. The conventional account places the symbol
as primary and the motor act as secondary: the typist knows the letter and
retrieves the corresponding key position. The attractor landscape account
inverts this: the motor landscape is primary and the symbol is the name of
a region within it. Expertise is not the acquisition of a larger symbol
inventory but the progressive refinement of a motor landscape whose stable
regions become better separated, more robustly attracted, and more richly
connected by learned trajectories.

Writing systems, on this view, are not inventories stored in memory. They are
navigable motor manifolds. Letters are stable regions within those manifolds.
Expertise progressively collapses symbolic retrieval into trajectory navigation,
until words become places and sentences become paths through a learned motor
landscape. The Eight-Letter Keyboard is a design that takes this seriously from
the beginning, organizing the motor landscape around the natural dichotomies of
the hand rather than the historical accidents of spatial key arrangement.

\section{The Motor Manifold: A Formal Treatment}

The arguments of this paper can be given a more precise mathematical form that
reveals their connection to a broader theory of symbolic projection. This
section develops that form, drawing on the same projection geometry that
underlies the admissibility framework in theories of representation and
reachability.

\subsection{Symbols as Quotient Classes of Motor Space}

Let $\mathcal{M}$ denote the space of possible hand configurations and
trajectories: the motor manifold. A motor act is a point $m \in \mathcal{M}$,
and a keyboard layout defines a projection
\[
  \sigma : \mathcal{M} \;\to\; \Sigma,
\]
where $\Sigma$ is the symbol space. For a conventional Latin keyboard,
$\Sigma = \{A, B, C, \ldots, Z\}$ together with punctuation and modifier
symbols. The projection $\sigma$ is many-to-one: many distinct motor
trajectories all successfully produce the same symbol. The preimage
$\sigma^{-1}(s)$ for any symbol $s \in \Sigma$ is therefore a set of
motor acts rather than a single point.

\begin{definition}[Motor Equivalence]
Two motor acts $m_1, m_2 \in \mathcal{M}$ are motor-equivalent with
respect to $\sigma$ if $\sigma(m_1) = \sigma(m_2)$. Write
$m_1 \sim_\sigma m_2$.
\end{definition}

Under this equivalence, the symbol space is the quotient
\[
  \Sigma \;\cong\; \mathcal{M} / {\sim_\sigma}.
\]
The alphabet is not a primitive inventory. It is a quotient space of motor
behavior: a set of equivalence classes on the space of possible hand movements.
This observation relocates the locus of symbolic identity from an abstract
character set to the motor structure of the body that generates it. Letters
are not given; they are carved from motor space by the projection $\sigma$.

\subsection{Generative Sufficiency}

The Eight-Letter Keyboard proposes a specific factored structure for this
projection. Rather than treating each letter as an independent equivalence
class with no internal structure, it proposes that the symbol space factors
as a product of simpler categorical dimensions.

\begin{definition}[Factored Motor Projection]
A projection $\sigma : \mathcal{M} \to \Sigma$ is \emph{factored} with
respect to feature sets $F$ (finger identity) and $D$ (directional parameter)
if there exists an encoding $\phi : F \times D \to \Sigma$ such that
$\sigma = \phi \circ (\pi_F, \pi_D)$, where $\pi_F$ and $\pi_D$ are the
projections of each motor act onto its finger-identity and directional
components respectively.
\end{definition}

\begin{theorem}[Generative Sufficiency]
If $|F| \cdot |D| \geq |\Sigma|$, then a surjective encoding
$\phi : F \times D \to \Sigma$ exists. In particular, the Latin alphabet
of twenty-six letters is recoverable from eight finger identities and three
directional positions per finger, since $8 \times 3 = 24 < 26$ requires only
that the index fingers each supply one additional position from a second column,
giving $6 + 6 + 3 + 3 + 3 + 3 = 24 + 6 = 30 \geq 26$.
\end{theorem}

\begin{proof}
Any surjection from a finite set of cardinality $n$ to a finite set of
cardinality $k \leq n$ exists by elementary combinatorics. The available
motor coordinate space $F \times D$, augmented by the two-column structure
of the index fingers, has cardinality $30$. Since $30 \geq 26$, a surjective
assignment of all twenty-six letters onto motor coordinates exists.
\end{proof}

This theorem is the formal justification for the Eight-Letter Keyboard. The
alphabet fits inside the finger-direction product space; no additional physical
resources are required beyond the combinatorial structure of the eight fingers
and their directional degrees of freedom.

\subsection{Motor Phonological Equivalence}

The analogy between typing and speech that motivates the earlier sections of
this paper can now be stated as a structural equivalence between two projection
systems.

\begin{theorem}[Motor-Phonological Equivalence]
Any symbolic system generated by a finite set of motor feature dimensions
$M_1, M_2, \ldots, M_k$ with respective cardinalities $n_1, n_2, \ldots, n_k$
is structurally equivalent, as a combinatorial symbol-generation system, to
a finite feature phonology over the same dimensions.
\end{theorem}

\begin{proof}
Both systems generate a symbol space as a quotient of a product space:
$\Sigma \cong \prod_i A_i / {\sim}$ for some equivalence relation $\sim$
capturing the many-to-one character of the projection. Whether the factors
$A_i$ are articulatory dimensions (place, manner, voicing) or motor dimensions
(finger identity, directional parameter, force class) is irrelevant to the
algebraic structure. The generation procedure is identical in both cases:
select a value from each dimension, and the combination determines a symbolic
output up to the equivalence induced by the projection. The algebra of finite
feature systems is the same whether instantiated in vocal or manual motor space.
\end{proof}

The practical implication is that the entire apparatus of phonological theory---
distinctive features, markedness, feature hierarchies, underspecification,
assimilation, and neutralization---is in principle available for the analysis
of keyboard motor systems. Keyboard layouts are a kind of manual phonology,
and the principles that govern the learnability, stability, and expressive
range of phonological systems can be brought to bear on the design of input
interfaces.

\subsection{Script Projection and Motor Preservation}

The script-projection keyboards discussed in Section~3 can be given a compact
formal treatment. Let $\mathcal{Q}$ denote the motor program acquired through
QWERTY training: the set of motor acts and their associated categorical
structures that constitute the QWERTY motor phonology. A script-projection
keyboard defines a mapping
\[
  f : \mathcal{Q} \;\to\; S,
\]
where $S$ is the target script (Arabic, Phoenician, Standard Galactic, or
Unicode Italic). The total learning cost $C$ of acquiring the new script via
this keyboard decomposes as
\[
  C \;=\; C_{\mathcal{M}} + C_S,
\]
where $C_{\mathcal{M}}$ is the motor reorganization cost and $C_S$ is the
symbolic remapping cost.

\begin{theorem}[Motor Preservation]
If $f$ is a bijection on motor channels---that is, if each key in $\mathcal{Q}$
maps to exactly one symbol in $S$ and the motor acts required to produce that
symbol are identical to those already learned---then $C_{\mathcal{M}} = 0$
and the total learning cost reduces to $C = C_S$.
\end{theorem}

\begin{proof}
Motor preservation means that no new motor programs need to be acquired or
old ones modified. The motor component of $C$ is therefore zero. The remaining
cost is the cost of learning the new symbol-to-glyph mapping, which is a
purely cognitive rather than motor task. \end{proof}

This theorem formalizes the intuition behind the script-projection design
principle. By holding the motor program fixed and varying only the symbolic
output, the projection keyboards reduce a potentially large acquisition burden
to the cost of learning a new interpretation of familiar gestures. The Arabic
and Phoenician keyboards further reduce $C_S$ by embedding phonetic and
etymological mnemonics into the projection, creating bridges between the new
symbolic layer and existing linguistic knowledge.

A broader connection to embodied cognition is worth making explicit here.
Contemporary gesture recognition research has converged on a structurally
parallel observation in a very different domain: the systems that perform
best at recognizing sign language, swipe-based text input, and musical gesture
are those that treat observable signals not as direct symbolic content but as
partial projections of a latent motor state, and reconstruct the intended
meaning through constraint-compatible inference over learned priors about
admissible motor trajectories. The observable gesture is evidence, not the
meaning. The meaning is recovered by finding the latent trajectory of highest
plausibility consistent with the evidence.

This is the inverse of the keyboard design problem, but the underlying geometry
is the same. The keyboard designer asks: given a learned motor vocabulary, what
projection makes the target symbolic domain most reachable? The gesture
recognition system asks: given an observed projection, what latent motor
intention best explains it? Both questions are about the relationship between
a motor manifold and a symbolic manifold, and both require the same concept
of admissibility: the set of gesture-projection pairs that are coherent, stable,
and semantically interpretable.

The mimetic cognition account of musical perception makes the connection vivid
from yet another direction. Listeners reconstruct the intended embodied action
of a musical performer from partial acoustic traces, implicitly simulating the
bodily gestures capable of producing the observed sound. This reconstruction is
not feature-matching but inverse inference: the listener models the admissibility
landscape of the performer's motor system and finds the most plausible trajectory
through it consistent with what is heard. The same mechanism that allows a
listener to infer a violinist's bow pressure from an acoustic envelope is, at
the appropriate level of abstraction, the same mechanism that allows a swipe
keyboard to infer an intended word from a finger trajectory, and the same
mechanism that allows a skilled typist to produce a word from a canonical
gesture without consciously planning the individual keystrokes. All three are
instances of constraint-compatible reconstruction over a learned motor
admissibility landscape.

\subsection{The Ark Principle Formalized}

The generative preservation argument of Section~5 can be stated as a formal
theorem connecting the keyboard design problems to the broader question of
symbolic survival.

\begin{definition}[Generator Sufficiency]
A set $G$ is a \emph{sufficient generator} for a symbolic system $S$ under
reconstruction operator $\rho$ if $\rho(G) = S$: every element of $S$ can
be reconstructed from $G$ by applying $\rho$.
\end{definition}

\begin{theorem}[Ark Principle]
Preservation of a symbolic system $S$ through a bottleneck does not require
carrying all of $S$. It requires carrying a sufficient generator $G$ together
with the reconstruction operator $\rho$.
\end{theorem}

\begin{proof}
If $G$ is carried through the bottleneck and $\rho$ is available on the far
side, then $\rho(G) = S$ guarantees that every element of $S$ is recoverable.
No additional information is required, because the generative structure
encoded in $G$ and $\rho$ is preservation-equivalent to the full inventory $S$.
\end{proof}

For the keyboard family discussed in this paper, the generator is the motor
phonology of the hand---the eight finger identities and their directional
degrees of freedom---and the reconstruction operator is the projection
$\phi : F \times D \to \Sigma$ that assigns symbols to motor coordinates.
The Latin alphabet, Arabic script, Phoenician alphabet, and Standard Galactic
glyphs are all recoverable from this generator once the appropriate projection
has been specified. The motor program of the hand is, in this precise sense,
an ark for symbolic civilization: it carries enough structure to reach an
enormous space of distinction systems from a single generative substrate.

This connects the keyboard design problem to a recurring pattern in the
theory of complex systems. What survives a bottleneck is never the full
inventory. What survives is the generative structure---the seeds, the grammar,
the motor program, the genetic code---from which the inventory can be
reconstructed. The question that good design in this space must answer is
therefore not how to carry more but how to identify the sufficient generator
that carries everything essential in the smallest possible form.

\bigskip

\noindent\textbf{Acknowledgements.} The hands knew before the theory caught up.

\begin{thebibliography}{99}

\bibitem{salthouse1986}
Salthouse, Timothy~A.
``Perceptual, Cognitive, and Motoric Aspects of Transcription Typing.''
\textit{Psychological Bulletin}, 99(3):303--319, 1986.

\bibitem{rumelhart1982}
Rumelhart, David~E. and Norman, Donald~A.
``Simulating a Skilled Typist: A Study of Skilled Cognitive-Motor Performance.''
\textit{Cognitive Science}, 6(1):1--36, 1982.

\bibitem{verwey1996}
Verwey, Willem~B.
``Buffer Loading and Chunking in Typewriting.''
\textit{Journal of Experimental Psychology: Human Perception and Performance},
22(2):544--562, 1996.

\bibitem{liberman1957}
Liberman, Alvin~M., Harris, Katherine~S., Hoffman, Howard~S., and Griffith,
Belinda~C.
``The Discrimination of Speech Sounds Within and Across Phoneme Boundaries.''
\textit{Journal of Experimental Psychology}, 54(5):358--368, 1957.

\bibitem{eimas1971}
Eimas, Peter~D., Siqueland, Einar~R., Jusczyk, Peter, and Vigorito, James.
``Speech Perception in Infants.''
\textit{Science}, 171(3968):303--306, 1971.

\bibitem{ladefoged2011}
Ladefoged, Peter and Johnson, Keith.
\textit{A Course in Phonetics}, 6th ed.
Wadsworth, 2011.

\bibitem{lashley1951}
Lashley, Karl~S.
``The Problem of Serial Order in Behavior.''
In Jeffress (ed.), \textit{Cerebral Mechanisms in Behavior}.
Wiley, 1951.

\bibitem{hikosaka1999}
Hikosaka, Okihide, Nakamura, Kae, Sakai, Katsuyuki, and Nakahara, Hiroyuki.
``Neural Correlates of Motor Learning, Memory, and Representation.''
\textit{Current Opinion in Neurobiology}, 9(2):217--227, 1999.

\bibitem{soechting1992}
Soechting, John~F. and Flanders, Martha.
``Moving in Three-Dimensional Space: Frames of Reference, Vectors, and
Coordinate Systems.''
\textit{Annual Review of Neuroscience}, 15:167--191, 1992.

\bibitem{clark1998}
Clark, Andy and Chalmers, David~J.
``The Extended Mind.''
\textit{Analysis}, 58(1):7--19, 1998.

\bibitem{hutchins1995}
Hutchins, Edwin.
\textit{Cognition in the Wild}.
MIT Press, 1995.

\bibitem{varela1991}
Varela, Francisco~J., Thompson, Evan, and Rosch, Eleanor.
\textit{The Embodied Mind}.
MIT Press, 1991.

\bibitem{gibson1979}
Gibson, James~J.
\textit{The Ecological Approach to Visual Perception}.
Houghton Mifflin, 1979.

\bibitem{miller1956}
Miller, George~A.
``The Magical Number Seven, Plus or Minus Two: Some Limits on Our Capacity
for Processing Information.''
\textit{Psychological Review}, 63(2):81--97, 1956.

\bibitem{chomsky1968}
Chomsky, Noam and Halle, Morris.
\textit{The Sound Pattern of English}.
Harper and Row, 1968.

\bibitem{kolmogorov1965}
Kolmogorov, Andrey~N.
``Three Approaches to the Quantitative Definition of Information.''
\textit{Problems of Information Transmission}, 1(1):1--7, 1965.

\bibitem{ashby1956}
Ashby, W.~Ross.
\textit{An Introduction to Cybernetics}.
Chapman \& Hall, 1956.

\bibitem{dehaene2005}
Dehaene, Stanislas.
\textit{Reading in the Brain}.
Viking, 2009.

\bibitem{elman1990}
Elman, Jeffrey~L.
``Finding Structure in Time.''
\textit{Cognitive Science}, 14(2):179--211, 1990.

\bibitem{iverson2005}
Iverson, Jana~M. and Goldin-Meadow, Susan.
``Gesture Paves the Way for Language Development.''
\textit{Psychological Science}, 16(5):367--371, 2005.

\bibitem{goldinmeadow2003}
Goldin-Meadow, Susan.
\textit{Hearing Gesture: How Our Hands Help Us Think}.
Harvard University Press, 2003.

\bibitem{cox2016}
Cox, Arnie.
\textit{Music and Embodied Cognition: Listening, Moving, Feeling, and Thinking}.
Indiana University Press, 2016.

\bibitem{zhai2012}
Zhai, Shumin and Kristensson, Per~Ola.
``The Word-Gesture Keyboard: Reimagining Keyboard Interaction.''
\textit{Communications of the ACM}, 55(9), 2012.

\bibitem{schieber2001}
Schieber, Marc~H.
``Individuated Finger Movements: Limb Kinematics, Physiology, and Motor
Control.''
\textit{Journal of Neurophysiology}, 2001.

\end{thebibliography}

\end{document}
